\chapter{Cleft Lip and Palate Repair}

\section*{Introduction \& Clinical Indications}
Cleft lip and cleft palate are congenital openings caused by incomplete fusion of the developing facial processes. Repair is planned by a multidisciplinary cleft or craniofacial team and may involve several operations as the child grows. Primary lip repair is commonly performed during the first year of life and palate repair during the first 18 months, but timing depends on the cleft, airway and feeding status, associated conditions, and the team's protocol.

The goals are to restore separation between the oral and nasal cavities, support feeding and speech development, improve lip and nasal form, and preserve facial growth. Additional procedures may later be required for fistula, velopharyngeal dysfunction, alveolar bone defects, orthodontic needs, or residual nasal or facial differences.

\section*{Relevant Surgical Anatomy}
The lip repair uses the philtral columns, vermilion border, Cupid's bow, nasal sill, and orbicularis oris muscle as alignment landmarks. Palatoplasty reconstructs the oral and nasal mucosa and reorients the muscles of the soft palate, especially the levator veli palatini, to support velopharyngeal function. The infraorbital and greater palatine neurovascular structures, facial vessels, tooth buds, and developing maxilla require protection. The exact anatomy varies with laterality, width, alveolar involvement, and syndromic features.

\section*{Preoperative Workup \& Preparation}
Assessment includes feeding and growth, airway and sleep symptoms, hearing and otitis history, dental and maxillary development, speech and language, and examination for associated anomalies or a syndrome. The team may include plastic or craniofacial surgery, pediatrics, anesthesia, otolaryngology, audiology, speech-language pathology, dentistry or orthodontics, genetics, nutrition, and psychosocial services. Routine imaging and blood tests are selective rather than universal; imaging is chosen when it will change operative planning.

The child must be medically optimized for anesthesia, with age-appropriate fasting instructions and a plan for analgesia and feeding after surgery. Consent should address bleeding, infection, airway events, wound or palatal fistula, altered facial growth, speech or resonance problems, dental and orthodontic treatment, and possible revision surgery. Active respiratory or systemic infection may require postponement.

\section*{Surgical Technique \& Procedure}
\begin{enumerate}
  \item \textbf{Positioning and anesthesia}: General anesthesia with a secured airway is used, with careful protection of pressure points and the eyes.
  \item \textbf{Lip repair}: Incisions are designed around the cleft margins. The orbicularis oris is released and reapproximated, the nasal floor and sill are reconstructed as needed, and the skin and vermilion are aligned in layers.
  \item \textbf{Palate repair}: The cleft edges are raised as mucoperiosteal or mucosal flaps. The nasal layer, palatal muscle sling, and oral mucosa are closed in layers without undue tension. The technique is selected for the cleft anatomy and surgeon's protocol.
  \item \textbf{Completion}: Hemostasis and airway safety are confirmed. The specimen is not ordinarily sent for pathology unless another lesion is present.
\end{enumerate}

\section*{Complications \& Risks}
Risks include bleeding, infection, airway obstruction, aspiration, wound separation, palatal fistula, dehiscence, scar, dental or maxillary growth disturbance, recurrent otitis media or hearing problems, velopharyngeal insufficiency, speech or resonance difficulty, and the need for revision. Risk varies with the cleft, age, associated conditions, and operation.

\section*{Postoperative Care \& Recovery}
Recovery includes airway and bleeding observation, multimodal analgesia, hydration, and a feeding plan appropriate to the repair and the child's baseline ability. The team gives specific instructions about permitted utensils, oral hygiene, arm restraints, and protecting the repair; these instructions vary by procedure. Fever, persistent bleeding, breathing difficulty, dehydration, wound separation, or inability to feed requires urgent review. Long-term follow-up monitors speech, hearing, dental development, facial growth, psychosocial well-being, and the need for secondary surgery.

\section*{Outcomes \& Prognosis}
Most children achieve good healing, but lip appearance, speech, hearing, dental development, and facial growth require longitudinal assessment. Outcomes cannot be represented by one universal success rate. A coordinated team approach improves continuity of care and allows treatment to be adjusted as the child develops.

\section*{Additional Clinical Considerations}
Timing is coordinated rather than dictated by a single age rule. The team considers the child's weight, anemia, airway, cardiopulmonary status, feeding efficiency, and associated syndromes. Lip repair is often planned during the first year, while palatal repair is commonly completed during the first 18 months to support speech development; a child with significant airway disease, poor growth, or a complex syndrome may require a modified sequence. The decision should be documented in terms of the functional goal and the anatomy being treated.

Feeding support begins before surgery. Infants with a palatal cleft may require specialized bottles, paced feeds, upright positioning, and a dietitian's input because they cannot generate normal suction. Nasal regurgitation is common before repair, but coughing, recurrent chest infections, fatigue, prolonged feeds, or poor weight gain require assessment for aspiration or inadequate intake. After surgery, the team explains whether a soft or liquid diet is required, whether a particular bottle or spoon is allowed, and how to maintain hydration without traumatizing the repair. Caregivers should never improvise oral appliances or manipulate the incision without instruction.

Hearing and speech follow-up are part of the operation's outcome, not optional cosmetic aftercare. Children with cleft palate have a high risk of eustachian-tube dysfunction and otitis media with effusion. Audiology and otolaryngology review may be needed, and tympanostomy tubes are considered according to hearing findings and local pediatric guidance. Speech-language assessment monitors articulation, resonance, nasal airflow, and language development. Velopharyngeal insufficiency may persist despite an anatomically healed palate and can require speech therapy, imaging or endoscopy, and secondary surgery such as a pharyngeal flap or sphincter pharyngoplasty.

The alveolus and maxilla require a separate long-term plan. Children with an alveolar cleft may need orthodontic expansion followed by secondary alveolar bone grafting, often during mixed dentition when the canine is developing. The timing depends on dental imaging, eruption stage, cleft width, and the child's growth. Orthodontic treatment, dental hygiene, caries prevention, and monitoring of tooth eruption should be coordinated with the craniofacial team. A primary lip or palate repair does not close every alveolar defect and should not be represented as completing the entire cleft pathway.

Potential revision is discussed without implying failure of the original operation. Scar management, lip or nasal revision, fistula closure, velopharyngeal surgery, orthodontics, and orthognathic surgery may be appropriate at different stages. Psychological support and age-appropriate participation in decisions become increasingly important as the child grows. Photographs and operative records can help the team compare symmetry and function over time, but appearance should be assessed in the context of breathing, feeding, speech, hearing, and the child's own goals.

During follow-up, new bleeding, fever, respiratory distress, dehydration, wound separation, foul drainage, or a sudden change in speech or swallowing requires review. A small fistula may present later as nasal leakage, recurrent nasal congestion, or an unexplained speech change rather than as an obvious wound problem. Long-term surveillance continues through facial growth and dental maturity because the effects of the cleft and its repairs evolve over time.

\section*{Additional Follow-Up Considerations}
Care is staged across childhood. The team monitors feeding and growth in infancy, hearing and middle-ear disease during early childhood, speech and resonance during language development, dental eruption and maxillary growth during school years, and psychosocial well-being throughout. A child who has healed well at the lip may still need treatment of the palate, alveolus, hearing, or speech. Families should be given a clear schedule and a contact point for concerns between appointments.

Operations are planned around function and growth. Secondary lip or nasal revision, fistula closure, pharyngeal surgery, alveolar bone grafting, orthodontics, and orthognathic surgery are separate decisions, not automatic stages for every child. The craniofacial team should discuss expected benefits and alternatives, especially when a procedure may affect facial growth, speech, dental development, or the airway. Long-term photographs and speech or hearing assessments can help measure progress without reducing the outcome to appearance alone.

\section*{Team-Based Long-Term Care}
The cleft team should maintain a longitudinal record of operations, hearing tests, speech evaluations, dental findings, growth, and psychosocial needs. This prevents a child from being lost between surgical stages and helps the team distinguish an expected developmental change from a complication. Primary-care clinicians and school services may need guidance about hearing, speech, feeding, and the timing of dental or orthodontic referrals.

Parents and older children should be involved in decisions at a level appropriate to age and maturity. The aims of revision should be explained in terms of function, appearance, or both, and the limits of surgery should be clear. A healed scar or a normal-looking lip does not guarantee normal speech, hearing, or maxillary growth. Conversely, a residual difference in appearance does not mean that the repair has failed if feeding, speech, airway, and psychosocial outcomes are satisfactory.

\section*{Patient and Family Education}
Families should receive written instructions in an accessible format, including feeding technique, pain medicine dosing, wound protection, emergency symptoms, and the date of the next team visit. Children should not be judged against a single photograph or speech sample; progress is measured over development. Families may benefit from peer support, social work, genetics counseling, and psychological care when appearance, speech, feeding, or repeated operations create distress.

\section*{Transition to Adult Care}
As adolescents approach adulthood, records and responsibilities should transition gradually to the patient. The team reviews residual speech, dental, hearing, nasal, facial, airway, and psychosocial concerns and identifies which issues need adult craniofacial, dental, ENT, or reproductive counseling. A completed operation list and imaging history reduce repeated testing and support informed decisions about any future revision.

\section*{Documentation and Safety}
The record should identify the cleft anatomy, technique, layers repaired, airway and feeding plan, complications, and planned team follow-up. Families should receive a written contact pathway and understand that bleeding, breathing difficulty, dehydration, fever, or wound separation is urgent. Longitudinal documentation allows speech, hearing, dental, and facial outcomes to be interpreted across multiple operations.

\section*{Practical Follow-Up}
The team reviews wound healing, feeding, hydration, breathing, and pain soon after surgery. Later visits address hearing, speech, dental eruption, facial growth, and psychosocial health. A normal-looking incision does not eliminate the need for longitudinal cleft-team care.

\section*{Final Safety Note}
The timing and sequence of repair are individualized by the cleft team. Families should receive clear instructions for feeding, wound protection, warning symptoms, and long-term speech, hearing, dental, and psychosocial follow-up.

\section*{Completion Checklist}
Before discharge, confirm the airway, hydration, analgesia, feeding method, wound instructions, caregiver understanding, and follow-up appointment. Before the episode is closed, confirm hearing, speech, dental, growth, and psychosocial referrals are documented.

\section*{Handoff}
The cleft-team plan should be shared with primary care and the family so that urgent symptoms and future developmental reviews are not missed.

\section*{Final Note}
Longitudinal team care is essential because speech, hearing, dental development, and facial growth evolve over time.

\section*{References}
\begin{enumerate}
  \item Centers for Disease Control and Prevention. Cleft Lip/Cleft Palate. Updated January 8, 2026.
  \item American Cleft Palate--Craniofacial Association. Parameters for Evaluation and Treatment of Patients with Cleft Lip/Palate or Other Craniofacial Anomalies. Revised edition.
  \item American Academy of Pediatrics. Health Supervision for Children With Cleft Lip and/or Cleft Palate.
\end{enumerate}

\clearpage
